\begingroup
\small
\setlength{\parskip}{2pt}
\setlength{\abovedisplayskip}{5pt}
\setlength{\belowdisplayskip}{5pt}
\begin{center}
{\Large\bfseries\color{azulPDMI}
Quantos triângulos se escondem num polígono?}\\[3pt]
{\small Resolução da tarefa de exploração -- aula simulada (50 minutos)}
\end{center}

\begin{tcolorbox}[
colback=azulClaro,colframe=azulPDMI!40,boxrule=.5pt,
arc=2mm,left=8pt,right=8pt,top=6pt,bottom=6pt]
\textbf{Ideia matemática central.}
A propriedade da soma dos ângulos internos de um triângulo
($180^\circ$) permite determinar a soma dos ângulos internos de um
polígono convexo, desde que este seja decomposto em triângulos que
preencham toda a sua região interior sem sobreposições.
A exploração dos exemplos conduz à expressão
\[
\boxed{S=(n-2)\times180^\circ,\qquad n\geq3,}
\]
cuja validade geral exige uma justificação matemática.
\end{tcolorbox}

\section*{Parte I -- À procura de uma estratégia}

\subsection*{1.1. Quadrilátero e pentágono}

Uma estratégia possível é traçar, a partir de um vértice, diagonais
que decomponham o polígono em triângulos. Os desenhos seguintes
ilustram essa estratégia.

\begin{center}
\begin{minipage}{.47\linewidth}
\centering
\begin{tikzpicture}[scale=.67,line join=round]
\coordinate (A) at (.3,.2);
\coordinate (B) at (4.3,.1);
\coordinate (C) at (4.7,2.6);
\coordinate (D) at (.5,2.9);
\draw[thick,draw=azulPDMI,fill=azulClaro]
  (A)--(B)--(C)--(D)--cycle;
\draw[very thick,orange!85!black] (A)--(C);
\node[below left] at (A) {$A$};
\node[below right] at (B) {$B$};
\node[above right] at (C) {$C$};
\node[above left] at (D) {$D$};
\end{tikzpicture}\\[3pt]
{\small Quadrilátero: 2 triângulos}
\end{minipage}\hfill
\begin{minipage}{.47\linewidth}
\centering
\begin{tikzpicture}[scale=.64,line join=round]
\coordinate (A) at (.4,.2);
\coordinate (B) at (3.8,.1);
\coordinate (C) at (5,2.2);
\coordinate (D) at (2.8,3.4);
\coordinate (E) at (.1,2);
\draw[thick,draw=azulPDMI,fill=azulClaro]
  (A)--(B)--(C)--(D)--(E)--cycle;
\draw[very thick,orange!85!black] (A)--(C) (A)--(D);
\node[below left] at (A) {$A$};
\node[below right] at (B) {$B$};
\node[right] at (C) {$C$};
\node[above] at (D) {$D$};
\node[left] at (E) {$E$};
\end{tikzpicture}\\[3pt]
{\small Pentágono: 3 triângulos}
\end{minipage}
\end{center}

No quadrilátero $ABCD$, a diagonal $AC$ forma os triângulos
$ABC$ e $ACD$. Portanto,
\[
S_4=2\times180^\circ=\boxed{360^\circ}.
\]

No pentágono $ABCDE$, as diagonais $AC$ e $AD$ formam os
triângulos $ABC$, $ACD$ e $ADE$. Assim,
\[
S_5=3\times180^\circ=\boxed{540^\circ}.
\]

\textbf{Nota:} Esta é uma resolução possível, não a única.
Podem escolher-se outros vértices e outras diagonais, desde que se
obtenha uma decomposição válida. Um procedimento diferente, como
ligar um ponto interior aos vértices, exige ter em conta os ângulos
adicionais em torno desse ponto.

\subsection*{1.2. Comparação com a estratégia de um colega}

\textbf{Exemplo de resposta possível:}
«Escolhemos diagonais diferentes, mas obtivemos os mesmos resultados:
$360^\circ$ no quadrilátero e $540^\circ$ no pentágono.
Em ambos os casos, dividimos a figura em triângulos e utilizámos a
soma de $180^\circ$ em cada triângulo.»

A resposta é aberta: os alunos podem ter utilizado estratégias distintas
ou identificado dúvidas diferentes. Importa comparar os procedimentos,
explicar por que são válidos e não apenas verificar se os resultados
numéricos coincidem.

\Needspace{9\baselineskip}
\subsection*{Uma estratégia para experimentar -- alíneas a), b) e c)}

\textbf{a)} A diagonal $AC$ do quadrilátero origina dois triângulos,
$ABC$ e $ACD$.

\textbf{b)} No pentágono, escolhendo $A$ como vértice de partida,
traçam-se as diagonais $AC$ e $AD$, obtendo-se três triângulos:
$ABC$, $ACD$ e $ADE$.

\textbf{c)} Cada triângulo tem soma das amplitudes dos ângulos
internos igual a $180^\circ$. Logo,
\[
\begin{aligned}
S_4&=2\times180^\circ=360^\circ,\\
S_5&=3\times180^\circ=540^\circ.
\end{aligned}
\]

\subsection*{1.3. Justificação da soma dos ângulos dos triângulos}

As diagonais traçadas no interior dos polígonos convexos
originam triângulos cujas regiões interiores não se sobrepõem
e cuja reunião é exatamente a região interior de cada polígono.

Além disso, \textbf{os ângulos dos triângulos que têm o mesmo vértice
do polígono recompõem, sem lacunas nem sobreposições, o ângulo
interno desse vértice}. Por exemplo, no quadrilátero,
\[
\widehat{DAB}=\widehat{DAC}+\widehat{CAB}
\quad\text{e}\quad
\widehat{BCD}=\widehat{BCA}+\widehat{ACD}.
\]
Nos vértices $B$ e $D$, os ângulos dos triângulos coincidem
com os ângulos internos correspondentes do quadrilátero.

Assim, a soma de todos os ângulos internos dos triângulos é
\textbf{exatamente} a soma dos ângulos internos do polígono;
não são acrescentados ângulos novos em pontos interiores, porque
os segmentos utilizados têm extremidades nos vértices do polígono.

\section*{Parte II -- Da exploração à regra geral}

\subsection*{2.1. Decomposição e preenchimento da tabela}

Traçando as diagonais a partir de um mesmo vértice, obtém-se:

\begin{center}
\renewcommand{\arraystretch}{1.24}
\begin{tabularx}{\linewidth}{|
>{\raggedright\arraybackslash}X|
>{\centering\arraybackslash}p{2.0cm}|
>{\centering\arraybackslash}p{2.5cm}|
>{\centering\arraybackslash}p{3.25cm}|}
\hline
\rowcolor{azulClaro}
\textbf{Polígono} &
\textbf{N.º de lados} &
\textbf{N.º de triângulos} &
\textbf{Soma dos ângulos internos} \\\hline
Triângulo & 3 & 1 & $1\times180^\circ=180^\circ$ \\\hline
Quadrilátero & 4 & 2 & $2\times180^\circ=360^\circ$ \\\hline
Pentágono & 5 & 3 & $3\times180^\circ=540^\circ$ \\\hline
Hexágono & 6 & 4 & $4\times180^\circ=720^\circ$ \\\hline
\end{tabularx}
\end{center}

\textbf{Construções esperadas:} no triângulo não é preciso
traçar diagonais; no quadrilátero traça-se uma diagonal;
no pentágono, duas; no hexágono, três. Em todos os casos,
as diagonais partem de um único vértice e produzem,
respetivamente, $1$, $2$, $3$ e $4$ triângulos.

\subsection*{2.2. Relação entre o número de lados e o número de triângulos}

\textbf{Resposta possível:}
«O número de triângulos obtidos é sempre inferior em duas unidades
ao número de lados do polígono. Quando passamos de um polígono
de $n$ lados para outro de $n+1$ lados, obtemos mais um triângulo
pela mesma estratégia de decomposição.»

Se $m$ designar o número de triângulos, a relação é
\[
\boxed{m=n-2.}
\]
Nesta fase, os exemplos permitem \emph{conjeturar} a relação.
A questão 2.4 destina-se a justificar que ela é válida em geral.

\subsection*{2.3. Expressão para a soma dos ângulos internos}

Se um polígono convexo com $n$ lados se decompõe em
$n-2$ triângulos, e cada triângulo contribui com $180^\circ$,
a soma $S$ das amplitudes dos ângulos internos é
\[
\boxed{S=(n-2)\times180^\circ.}
\]

Na expressão, $n$ representa o número de lados,
$n-2$ o número de triângulos da decomposição,
e $180^\circ$ a soma das amplitudes dos ângulos
internos de cada triângulo. Por distributividade,
pode também escrever-se
\[
S=180^\circ n-360^\circ.
\]
As duas expressões são equivalentes.

\subsection*{2.4. Justificação para qualquer polígono convexo}

Considere-se um polígono convexo com $n$ lados,
sendo $n$ um número inteiro e $n\geq3$.
Escolha-se um dos seus vértices, por exemplo $A$.

Dos restantes $n-1$ vértices, dois são adjacentes a $A$:
os extremos dos dois lados que aí se encontram.
Os outros $n-3$ são não adjacentes, pelo que é possível
ligar $A$ a cada um deles através de uma diagonal.

Por o polígono ser \textbf{convexo}, essas diagonais
ficam no interior do polígono. Traçadas a partir do mesmo
vértice, não se cruzam no interior e dividem-no numa sequência
de regiões triangulares, sem sobreposições nem espaços por preencher.

Os $n-3$ segmentos interiores, juntamente com os dois lados
incidentes em $A$, originam $n-2$ setores triangulares:
\[
\underbrace{n-3}_{\text{diagonais}}
+\underbrace{1}_{\text{região triangular adicional}}
=\underbrace{n-2}_{\text{triângulos}}.
\]
Em particular, para $n=3$, não se traçam diagonais e
obtém-se o próprio triângulo, em conformidade com a expressão.

Como a soma das amplitudes dos ângulos internos dos triângulos
coincide com a soma das amplitudes dos ângulos internos do
polígono (questão 1.3), conclui-se que
\[
S=\underbrace{180^\circ+\cdots+180^\circ}_{n-2\text{ parcelas}}
=\boxed{(n-2)\times180^\circ}.
\]

\textbf{A justificação não depende de verificar alguns exemplos:}
explica por que razão a decomposição e a contagem dos triângulos
funcionam para qualquer polígono convexo, com $n\geq3$ lados.

\section*{Para concluir}

Uma forma de completar a frase da ficha é:

\begin{tcolorbox}[
colback=azulClaro,colframe=azulPDMI!40,boxrule=.5pt,
arc=2mm,left=8pt,right=8pt,top=6pt,bottom=6pt]
«Num polígono convexo com $n$ lados, a soma das amplitudes
dos ângulos internos é
\[
\boxed{S=(n-2)\times180^\circ}
\]
porque podemos dividir esse polígono em $n-2$ triângulos,
sem sobreposições nem espaços por preencher,
e a soma das amplitudes dos ângulos internos de cada
triângulo é $180^\circ$.»
\end{tcolorbox}

\section*{Se houver tempo -- polígono convexo com 12 lados}

Como $n=12$, a decomposição a partir de um vértice
origina $12-2=10$ triângulos. Logo,
\[
S=(12-2)\times180^\circ
=10\times180^\circ
=\boxed{1800^\circ}.
\]
Não é necessário desenhar o polígono, pois a justificação
da expressão geral garante a sua validade para $n=12$.

\endgroup
